\section{Composition, gluing and the spatial-kernel limitation}
\label{sec:pairings}

\subsection{Transport and fusion}

Write $\Phi(r)=[f^{(0)}(r),f^{(1)}(r)]$ and
\begin{equation}
 E_t=J_x^hJ_y^h b_t^{-2h}\Phi(a_t/b_t).
\end{equation}
In a single original-coordinate frame, the exact propagator is
\begin{equation}
 U(t,s)=E_tE_s^{-1},\qquad
 U(t,s)U(s,r)=U(t,r),\qquad U(s,s)=I.
 \label{eq:composition}
\end{equation}
It is also the time-ordered exponential of the generator. Restarting the
Loewner map requires the translated remaining driver and the evolved marked
points, rather than a reset to the original geometry. The primary-Jacobian
chain rule gives the same composition.

For fusion, define
\begin{equation}
 R=\begin{pmatrix}1/2&\sqrt3/2\\-\sqrt3/2&1/2\end{pmatrix},\quad
 Z=\diag(1,-1),\qquad \Phi_{12}(r)=RZ\Phi(1-r)Z.
\end{equation}
Its columns have unit leading tensors in the $12$ channels.
Direct substitution of \eqref{eq:block0}--\eqref{eq:block1} yields
\begin{equation}
 \Phi(r)=\Phi_{12}(r)\mathsf C,\qquad
 \mathsf C=\begin{pmatrix}
 1/\sqrt2&-\sqrt{2/3}\\\sqrt{3/8}&1/\sqrt2
 \end{pmatrix}.
\end{equation}
With $N=\diag(1,\sqrt3/2)$, the unitary channel gauge gives
\begin{equation}
 N^{-1}\mathsf C N=\frac1{\sqrt2}
 \begin{pmatrix}1&-1\\1&1\end{pmatrix}.
 \label{eq:unitary_fusion}
\end{equation}
The signs depend on the chosen tensor basis and channel gauge.

\subsection{The channel metric is not the fixed tensor metric}

In raw Frobenius coefficients $c,d$, the positive unitary Chern--Simons
pairing, normalized by vacuum norm one, is
\begin{equation}
 \ip{c}{d}_{\rm CS}=c^\dagger H_c d,
 \qquad H_c=\diag(1,4/3).
 \label{eq:channel_metric}
\end{equation}
Local monodromies have the distinct phases of the exponents $-3/8,1/8$,
forcing the invariant metric to be diagonal. Requiring fusion to preserve
it fixes the ratio $4/3$. Equivalently, \eqref{eq:unitary_fusion} is
unitary on $N^{-1}c$. Orientation reversal conjugates the amplitude and
sewing contracts the matching channel indices with this pairing. The
selected vacuum has norm one and equal-magnitude coefficients $1/\sqrt2$
in the normalized crossed channel.

For evaluated tensors $M=E_tc$, the same pairing becomes
\begin{equation}
 H_t=E_t^{-\dagger}H_cE_t^{-1},\qquad
 U(t,s)^\dagger H_tU(t,s)=H_s.
 \label{eq:transported_metric}
\end{equation}
Thus the CS channel norm is preserved while the evaluated tensor can
decrease in the fixed $W$ norm of \eqref{eq:tensor_balance}. Neither
metric has been identified with the arithmetic $L^2$ norm.
Likewise a scalar polarization $v^\dagger M$ with fixed $v\in W$ is
different from a channel amplitude $d^\dagger H_cc$.

\subsection{Spatial superposition does not create the initial identity}

For disjoint compact boundary intervals $I_x<I_y$ away from the hull,
\begin{equation}
 M_t[f,g]=\int_{I_x}\int_{I_y}f(x)g(y)M_t(x,y)\dd y\dd x
\end{equation}
is a well-defined superposition for smooth compactly supported tests.
Its derivative is the corresponding integral of $G_t(x,y)M_t(x,y)$.
Since $G_t$ depends on the positions, this single smeared two-vector does
not satisfy a closed $2\times2$ equation. Retaining all positions gives
a multiplication evolution, and integrating \eqref{eq:tensor_balance}
against a positive spatial measure gives a direct-integral norm balance.
Arbitrary functions in that enlarged space have not been prepared as
physical boundary states.

There is a stronger initial-condition test. Fix a compact interval
$I\subset(0,\infty)$ away from the hull and a polarization $v\in W$.
Set $\mathcal K_t(x,y)=v^\dagger M_t(x,y)$ for $x<y$, and zero for the
opposite order. This chosen ordering does not derive arithmetic causality.

\begin{proposition}[The raw ordered kernel has a compact initial limit]
On a compact time interval where $g_t$ is regular on $I$, the integral
operator with kernel $\mathcal K_t$ is Hilbert--Schmidt. As $t\downarrow0$
it converges in Hilbert--Schmidt norm to its explicit compact operator at
$t=0$, and therefore cannot have the identity as its initial limit on
the infinite-dimensional space $L^2(I)$.
\end{proposition}
\begin{proof}
The worst diagonal singularity is the $12$ identity OPE, of exponent
$-2h=-3/8$. Uniform regularity of the coordinate map and spectator
Jacobians gives
\begin{equation}
 |\mathcal K_t(x,y)|\le C_T(1+|x-y|^{-3/8}).
\end{equation}
The square has exponent $-3/4$, integrable on $I\times I$. Pointwise
continuity off the diagonal and dominated convergence yield the stated
Hilbert--Schmidt limit. That limit is compact and cannot equal the identity.
\end{proof}

A finite scalar normalization or a regular change of time cannot alter
this explicit limit. The result does not exclude a new sewing construction
with descendants, a separately derived identity term, or a distinct
singular limiting regime. Compact operators can converge strongly to the
identity in other settings; the present obstruction also uses convergence
to the specified compact initial operator.
